% =====================================================================
%  CAPÍTULO 3 — METODOLOGIA E TECNOLOGIAS
% =====================================================================
\chapter{Metodologia e Tecnologias}
\label{cap:metodologia}

\section{Metodologia de Desenvolvimento}
\label{sec:metodologia}

O desenvolvimento seguiu uma abordagem \textbf{ágil e incremental},
adaptada às características de um projeto académico individual. Em vez de
um modelo em cascata rígido, adotou-se um ciclo iterativo inspirado em
\emph{Kanban}: o trabalho foi organizado em incrementos verticais, cada
um entregando uma fatia funcional ponta-a-ponta (da ingestão à geração),
o que permitiu validar continuamente o \emph{pipeline} completo em vez de
esperar pela conclusão de todos os módulos.

\subsection{Fases do projeto}

O projeto desenrolou-se em fases sobrepostas:

\begin{enumerate}[label=\textbf{F\arabic*.},leftmargin=*]
  \item \textbf{Investigação e estado da arte} — levantamento das
  soluções existentes (\Cref{cap:estado_arte}), das tecnologias de IA
  generativa e das ferramentas de código aberto por módulo.
  \item \textbf{Prototipagem do \emph{pipeline}} — implementação inicial
  dos quatro módulos (M1--M4) em modo \emph{local-first}, com SQLite e
  ficheiros locais, validando a viabilidade técnica end-to-end.
  \item \textbf{Consolidação do \emph{backend}} — introdução de
  autenticação, fila de tarefas assíncronas, validação de \emph{uploads},
  \emph{rate limiting}, \emph{logging} estruturado e testes automatizados.
  \item \textbf{Construção do \emph{frontend}} — desenvolvimento da
  aplicação React completa (páginas, gestão de estado,
  internacionalização, tema).
  \item \textbf{Migração para a nuvem} — evolução para uma arquitetura
  multi-inquilino com Supabase (autenticação, base de dados e
  armazenamento) e \emph{deployment} em Render e Vercel.
  \item \textbf{Validação e documentação} — testes, recolha de resultados
  e redação do presente relatório.
\end{enumerate}

\subsection{Gestão e controlo de versões}

Todo o código foi versionado com \textbf{Git}, com \emph{commits}
descritivos a marcar cada incremento funcional e cada correção. O
histórico de versões serviu simultaneamente de registo de progresso e de
mecanismo de segurança (possibilidade de reverter alterações). A
\Cref{fig:metodologia} ilustra o cronograma das fases.

\begin{figure}[H]
\centering
% \includegraphics[width=\linewidth]{cronograma.png}
\figph[5.5cm]{Cronograma do projeto (diagrama de Gantt) com as fases
F1--F6 ao longo do tempo. Sugestão: criar em ferramenta de gestão de
projeto ou em TikZ.}
\caption{Cronograma das fases de desenvolvimento.}
\label{fig:metodologia}
\end{figure}

\section{\emph{Stack} Tecnológica e Justificação}
\label{sec:stack}

A seleção tecnológica obedeceu a três critérios: \emph{(i)} maturidade e
adequação técnica; \emph{(ii)} custo nulo ou camada gratuita; e
\emph{(iii)} possibilidade de execução local para garantir privacidade.
A \Cref{tab:stack} resume a \emph{stack}; as subsecções seguintes
justificam as escolhas principais.

\begin{table}[H]
\centering
\caption{\emph{Stack} tecnológica por camada.}
\label{tab:stack}
\small
\begin{tabularx}{\linewidth}{@{}p{3cm}Y@{}}
\toprule
\textbf{Camada} & \textbf{Tecnologias} \\
\midrule
\emph{Frontend} & React~18, TypeScript (modo estrito), Vite~5, Tailwind
CSS~3, TanStack Query, Zustand, React Router~6, react-i18next, Framer
Motion, React Flow, Sonner, Lucide. \\
\emph{Backend} & Python~3, FastAPI (assíncrono), SQLAlchemy~2 (async),
Pydantic~v2, slowapi, structlog, Celery + Redis (opcional). \\
IA / ML & Ollama (Llama~3.2~3B, local); Google Gemini~1.5~Flash (nuvem);
ChromaDB + \emph{sentence-transformers} (RAG, opcional); Tesseract
(\textsc{ocr}); NetworkX (grafo). \\
Multimédia & MoviePy + FFmpeg; Pillow (Ken Burns / \emph{letterbox});
edge-tts / gTTS (\textsc{tts}). \\
Dados & Supabase Auth (JWT/JWKS), PostgreSQL (asyncpg), Supabase Storage;
SQLite (desenvolvimento local). \\
Operação & Docker; Render (\emph{backend}); Vercel (\emph{frontend});
Supabase (dados); Git. \\
\bottomrule
\end{tabularx}
\end{table}

\subsection{\emph{Backend}: FastAPI}

Escolheu-se o \textbf{FastAPI} pela sua natureza \emph{assíncrona}
(essencial para um sistema com operações de I/O e de IA potencialmente
longas), pela validação automática de dados via \textbf{Pydantic} (que
elimina uma classe inteira de erros de \emph{input}) e pela geração
automática de documentação interativa (\emph{Swagger UI}). A combinação
com \textbf{SQLAlchemy} assíncrono e \textbf{Pydantic~v2} oferece tipagem
forte de ponta a ponta no servidor.

\subsection{IA: estratégia híbrida local/nuvem}

A decisão mais estruturante foi a \textbf{estratégia de IA híbrida}.
Privilegia-se a inferência \textbf{local} via \textbf{Ollama} (executando
o \emph{Llama~3.2~3B}) por três razões: \emph{privacidade} (os factos do
arquivo não saem da máquina), \emph{custo} (nulo) e \emph{autonomia} (não
depende de conectividade). Como \emph{fallback}, recorre-se ao
\textbf{Google Gemini~1.5~Flash}, cuja camada gratuita é suficiente para
um protótipo académico e cuja maior capacidade é útil para narrativas
mais exigentes ou quando não há \textsc{gpu} local. Para a análise visual
das fotografias (M1) usa-se o \textbf{Gemini Vision}. Esta arquitetura
de dupla via é detalhada no \Cref{cap:implementacao}.

\subsection{RAG: ChromaDB}

Para o subsistema \textbf{RAG} adotou-se o \textbf{ChromaDB} como base
vetorial, com \emph{embeddings} de \emph{sentence-transformers}. A
escolha recaiu sobre o ChromaDB pela simplicidade de integração
(persistência local em ficheiro, sem servidor dedicado). Como se verá, o
subsistema foi desenhado para \emph{degradar graciosamente} na ausência
desta dependência, uma decisão motivada pelas restrições de
\emph{deployment} discutidas no \Cref{cap:implementacao}.

\subsection{Multimédia: MoviePy, FFmpeg e Pillow}

Para a Abordagem A (\Cref{sub:geracao_video}) usaram-se o \textbf{MoviePy}
(composição programática de vídeo), o \textbf{FFmpeg} (codificação e
mistura de áudio) e o \textbf{Pillow} (processamento de imagem para o
efeito \emph{Ken~Burns} e para o enquadramento \emph{letterbox}). Para a
narração, optou-se pelas vozes neuronais europeias do \textbf{edge-tts},
com o \textbf{gTTS} como alternativa.

\subsection{\emph{Frontend}: React, TypeScript e TanStack Query}

No \emph{frontend} escolheu-se o \textbf{React} com \textbf{TypeScript em
modo estrito} para capturar, em tempo de compilação, incompatibilidades
de esquema com o \emph{backend}. O \textbf{Vite} foi preferido como
\emph{bundler} pela velocidade do \emph{Hot Module Replacement}. Para a
gestão de estado, distinguiu-se rigorosamente o \emph{estado de servidor}
— gerido pelo \textbf{TanStack Query} (cache, invalidação,
\emph{polling}) — do \emph{estado de cliente} — gerido pelo
\textbf{Zustand} (autenticação e tema). Esta separação, hoje considerada
boa prática, evita a complexidade de soluções como o Redux para dados que
são, na sua maioria, \emph{server-state}.

\subsection{Dados e operação: Supabase, Docker, Render e Vercel}

A camada de dados foi externalizada para o \textbf{Supabase}, que fornece
de forma integrada autenticação (com tokens JWT validados por JWKS),
base de dados \textbf{PostgreSQL} e armazenamento de objetos
(\emph{Storage}). O \emph{backend} é empacotado em \textbf{Docker} (para
incluir dependências de sistema como FFmpeg, Tesseract e libmagic) e
implantado no \textbf{Render}; o \emph{frontend} é servido pela
\textbf{Vercel}. Esta combinação assenta integralmente em camadas
gratuitas, mantendo o custo operacional próximo de zero.
